% Results for Section 7 (the opportunity-security space), to be pasted into
% \subsection{Results}\label{subsec:joint_results} of sections/07_joint_space.tex.
% Generated numbers: paper/results/joint.json; figures and table: scripts/paper/fig_joint.R
% (final rerun of 2026-10-12: cells rankable on both pillars only, two-tier coalition of 57
% members, Strategic Index with the IEA PAMS-required policy basis and the endowment screen on
% the allied picks, upstream trade baskets for EV, Wind, Nuclear, Solar and the grid).
% Figures: joint_scatter, joint_corr, joint_quadrants, joint_members. Table: tables/joint_stage.tex.
% Label mapping: fig:joint_scatter and fig:joint_quadrants REPLACE the two placeholder figure
% environments of Section 7 (same labels); fig:joint_corr, fig:joint_members and tab:joint_stage
% are new. The five \tbd paragraphs of the Section 7 Results subsection are superseded by this
% fragment.

\paragraph{Sample.} Both pillars score 6{,}232 country cells in the eleven headline technologies,
for 219 economies. Cross-country comparisons use only the cells that are rankable on both pillars
(Section~\ref{subsec:norm_diagnostics}): the economy is sovereign and at most half of the
pillar's weight in the cell is imputed. This leaves 3{,}157 cells for 175 economies. Of the
3{,}075 cells set aside, 607 belong to territories and 2{,}468 to sovereign economies whose
imputed weight share exceeds one half in Energy Security only (721), in Economic Opportunity only
(1{,}006) or in both (741). The cells fall in 30 of the 33 technology-stage pairs; Energy
Security does not score downstream grid or upstream and midstream geothermal, and neither
pillar's diagnostics flag any remaining pair as constant or unscored. The number of rankable
economies per pair ranges from 40 in midstream green hydrogen to 175 in upstream nuclear. Both
pillars are built from ranks within technology-stage pairs, but their levels differ
systematically across pairs: technology-stage fixed effects account for 41.9 per cent of the
variance of $\ES_{its}$ and 53.4 per cent of the variance of $\EO_{its}$ in this sample. All
statements below are therefore made within pairs unless they concern a country's own sectors.
Correlations across economies are Spearman rank correlations within a pair, and the four regions
of the space are defined by the within-pair medians of the two pillars over the rankable cells, a
cell at the median counting as low. All 57 members of the coalition $\Coalition$ of
Section~\ref{sec:coalition}, 31 in the core and 26 partners, have rankable cells: 1{,}500 in all,
903 for the core and 597 for the partners. The six European Union members that enter the allied
network only through the EU node are not counted as members here.

\paragraph{Security and opportunity move together, increasingly so along the chain.} Across the
30 pairs, the median within-pair Spearman correlation between $\ES_{its}$ and $\EO_{its}$ is
0.59, with an interquartile range of 0.45 to 0.77 (Table~\ref{tab:joint_stage}). It exceeds 0.5
in 19 pairs and is negative in one, midstream green hydrogen ($-0.21$). The highest value is
0.94, in midstream electric vehicles. A Spearman correlation computed by pooling all cells on
pillar levels, 0.48, understates the within-pair association (0.57 when the within-pair
percentile ranks are pooled), because it mixes the association within pairs with differences in
pillar levels across pairs. The dominant pattern is a gradient by stage, visible in the three
panels of Figure~\ref{fig:joint_scatter}: the median correlation is 0.44 upstream, 0.58
midstream and 0.78 downstream. The gradient also appears within most technologies, so it is not
only a matter of which technologies are scored at which stage: for wind the correlation is 0.37
upstream, 0.72 midstream and 0.89 downstream, for nuclear 0.16, 0.49 and 0.89, for coal 0.35,
0.56 and 0.79, for gas 0.40, 0.54 and 0.78, and for oil 0.50, 0.55 and 0.75. Three technologies
depart from it. In solar (0.48, 0.78 and 0.58) and electric vehicles (0.60, 0.94 and 0.80) the
correlation peaks midstream. In green hydrogen the midstream correlation ($-0.21$) lies below the
upstream one (0.03), in a pair with only 40 rankable economies whose security score takes few
distinct values. In batteries the correlation falls along the chain, from 0.63 upstream to 0.61
midstream and 0.40 downstream. Across technologies, the median over a technology's pairs ranges
from 0.03 for green hydrogen to 0.80 for electric vehicles.

Upstream, the correlation is positive in all ten pairs, and its size follows what the two
pillars' trade measures have in common. The upstream trade baskets of the final data give the
opportunity pillar a Trade category in every upstream pair except green hydrogen. In the five
mineral pairs with a trade basket on both sides (batteries 0.63, grid 0.60, electric vehicles
0.60, solar 0.48 and wind 0.37), the opportunity side's Trade score and the security side's Minerals Trade score
measure exports of overlapping materials (the grid basket, for instance, consists of copper ores,
blister copper, anodes and cathodes, bauxite, alumina and unwrought aluminium), and their
within-pair rank correlations lie between 0.70 and 0.85. The fossil-fuel pairs follow, oil
(0.50), gas (0.40) and coal (0.35). The two lowest are nuclear (0.16), in which half of the
rankable economies share one value of $\ES_{its}$ (see below), and green hydrogen (0.03), the
only upstream pair without a trade basket: there security rests on the platinum-group and other
mineral endowments while opportunity draws on cost, demand, readiness and production components
that do not depend on them.

\begin{figure}[tbp]
  \centering
  \includegraphics[width=\textwidth]{joint_scatter}
  \caption{The opportunity-security space. Each point is a country cell $(i,t,s)$ that is
  rankable on both pillars (a sovereign economy with at most half of the pillar's weight
  imputed), plotted at its percentile rank across the rankable economies of its technology-stage
  pair on $\ES_{its}$ (horizontal) and on $\EO_{its}$ (vertical); ranks are mid-ranks scaled to
  the unit interval. Panels are supply-chain stages. Coalition members
  (Section~\ref{sec:coalition}) are shown in violet (the 31 core members) and pink (the 26
  partners), other economies in gray; the six EU members that belong only to the network's EU
  node count as other economies. The solid lines at rank 0.5 mark the within-pair medians and
  divide the four regions: opportunity without security (top left), leaders (top right),
  laggards (bottom left) and security without opportunity (bottom right). The percentage above or
  below each corner of a panel is the share of the stage's cells in that region, a cell at the
  median counting as low; the upstream panel names the regions, with ``opportunity only'' and
  ``security only'' standing for the two off-diagonal ones. Panel headers give the number of
  pairs, the median across them of the within-pair Spearman correlation between the two pillars,
  and the same median computed from the non-shared parts of the pillars, $B^{\ES}$ and $B^{\EO}$
  of equation~\eqref{eq:joint_cov} (renormalised); downstream the latter is defined in only three
  pairs, because in the other seven every component of $\ES_{its}$ is shared. Labels mark, in each
  stage, the two coalition cells with the largest excess of the opportunity rank over the security
  rank. Columns of points arise where many economies share one value of a pillar: the densest
  column in the upstream panel, just left of the security median, is upstream nuclear, in which 87 of the 175
  rankable economies, all without identified uranium resources or mine production, share the
  value 0.172 of $\ES_{its}$, which is also the pair's median. A tied cell is plotted at the mean
  rank of its tie group, so 29 cells (0.9 per cent), tied at the median of $\ES_{its}$ in
  downstream green hydrogen and batteries, appear at security ranks between 0.53 and 0.55, just
  to the right of the line, although they are counted on its low side. Sample: 3{,}157 cells, 175
  economies, 30 technology-stage pairs (Table~\ref{tab:joint_stage}). Source:
  \texttt{energy\_security\_index} and \texttt{economic\_opportunity\_index} in
  \texttt{index\_outputs.rds}; \texttt{scripts/paper/fig\_joint.R}.}
  \label{fig:joint_scatter}
\end{figure}

\paragraph{How much of the association is mechanical.} The stage gradient follows from what the
two pillars share (Section~\ref{subsec:eo_shared}, Table~\ref{tab:shared_signals}). We classify a
component as shared when the other pillar contains a component of the same category whose
observed values are identical. On this definition every component of Consumption, Energy Access,
Investment and Energy Prices is shared, as are two of the three Energy Security and two of the
five Economic Opportunity components of Trade (the world market share and the concentration of
world exports) and five of the six components of Foreign Dependency in each pillar; no Production
component is shared. Shared components carry on average 61 per cent of the effective weight of
$\ES_{its}$ but only 28 per cent of that of $\EO_{its}$, and their share in $\ES$ rises along the
chain, from 25 per cent upstream to 62 per cent midstream and 91 per cent downstream. In seven of
the ten downstream pairs (batteries, coal, electric vehicles, gas, green hydrogen, nuclear and
oil) every component of $\ES_{its}$ is a deployment or demand signal that also enters
$\EO_{its}$ with identical values: the security score in these pairs is a reweighting of part of
the opportunity score, which is why downstream green hydrogen (0.92) and nuclear (0.89) show some
of the highest correlations of all pairs. Recomputing both pillars from their non-shared
components only, which is possible in the remaining 23 pairs, lowers the median within-pair
correlation from 0.55 to 0.30 (interquartile range 0.20 to 0.40), a reduction of 45 per cent. In
the seven fully shared downstream pairs the co-movement is mechanical by construction.

The decomposition of equation~\eqref{eq:joint_cov}, applied to the pillar levels within each pair
and pooled over pairs, gives the same picture in cardinal terms (Table~\ref{tab:joint_stage}).
The level covariance is positive in all 30 pairs, including midstream green hydrogen, whose rank
correlation is negative. The shared term $\operatorname{Cov}(A^{\ES},A^{\EO})$ accounts for 61
per cent of the pooled within-pair covariance of $\ES_{its}$ and $\EO_{its}$, the two cross terms
for 24 per cent and the covariance of the non-shared parts $\operatorname{Cov}(B^{\ES},B^{\EO})$
for 16 per cent. The shared term's share is 7 per cent upstream, 42 per cent midstream and 87 per
cent downstream, where the non-shared term contributes nothing; upstream the non-shared term
carries 58 per cent. The pooled shares weight each pair by the size of its covariance, which is
largest in the downstream pairs; in the median pair the shared term accounts for 19 per cent. The
cross terms are not definitional: they measure how the inputs common to both pillars co-move with
the other pillar's own inputs. They do, however, run through inputs that the two pillars share.

What remains once the shared inputs are removed runs mainly through trade. The median non-shared
rank correlation is 0.37 midstream and 0.30 upstream. The pillars' distinctive trade components,
net imports relative to GDP and, upstream, the minerals-trade measures on the security side, and
revealed comparative advantage, export size and feasibility on the opportunity side, all reward
economies that export the stage's basket; removing the opportunity side's trade components as
well lowers the midstream median to 0.15 and the upstream median to 0.02. This link is not
definitional, since the components are different variables, but it reflects the weight that both
pillars give to export performance, and the upstream baskets of the final data have extended it
to the mineral pairs. In the three downstream pairs in which the non-shared correlation can be
computed (geothermal, solar and wind) it is $-0.04$. Net of shared inputs, security and
opportunity are close to unrelated across economies except where a country's exports reveal both.

% Generated by scripts/paper/fig_joint.R; do not edit by hand.
\begin{table}[tbp]
\centering
\begin{threeparttable}
\caption{Energy Security and Economic Opportunity within technology-stage pairs, by stage}
\label{tab:joint_stage}
\footnotesize
\begin{tabular}{@{}lcccc@{}}
\toprule
 & Upstream & Midstream & Downstream & All \\
\midrule
Technology-stage pairs & 10 & 10 & 10 & 30 \\
Country cells & 999 & 1,048 & 1,110 & 3,157 \\
\addlinespace
\multicolumn{5}{@{}l}{\textit{Spearman correlation of $\ES_{its}$ and $\EO_{its}$ within pairs}}\\
\quad Median across pairs & 0.44 & 0.58 & 0.78 & 0.59 \\
\quad Lowest pair & 0.03 & $-$0.21 & 0.40 & $-$0.21 \\
\quad Highest pair & 0.63 & 0.94 & 0.92 & 0.94 \\
\quad Pooled within-pair ranks & 0.39 & 0.60 & 0.71 & 0.57 \\
\addlinespace
\multicolumn{5}{@{}l}{\textit{Components shared by the two pillars}}\\
\quad Effective weight in $\ES$ & 0.25 & 0.62 & 0.91 & 0.61 \\
\quad Effective weight in $\EO$ & 0.26 & 0.20 & 0.37 & 0.28 \\
\quad Pairs with both non-shared pillars defined & 10 & 10 & 3 & 23 \\
\quad Median correlation, non-shared pillars & 0.30 & 0.37 & $-$0.04 & 0.30 \\
\quad Median correlation, full pillars, same pairs & 0.44 & 0.58 & 0.58 & 0.55 \\
\quad Share of level covariance, shared term & 0.07 & 0.42 & 0.87 & 0.61 \\
\quad Share of level covariance, non-shared term & 0.58 & 0.20 & 0.00 & 0.16 \\
\addlinespace
\multicolumn{5}{@{}l}{\textit{Share of cells by region of the space (percent)}}\\
\quad Opportunity without security & 19.5 & 13.8 & 12.4 & 15.1 \\
\quad Leaders (both high) & 30.2 & 36.0 & 37.2 & 34.6 \\
\quad Security without opportunity & 17.2 & 12.5 & 10.2 & 13.2 \\
\quad Laggards (both low) & 33.0 & 37.7 & 40.2 & 37.1 \\
\bottomrule
\end{tabular}
\begin{tablenotes}[flushleft]\footnotesize
\item \textit{Notes:} Eleven headline technologies; cells rankable on both pillars only (a sovereign economy with at most half of the pillar's weight imputed), which leaves 3,157 of the 6,232 cells that both pillars score, for 175 economies. A technology-stage pair enters when both pillars score it and neither pillar's diagnostics flag it as constant or unscored, which leaves 30 of the 33 pairs: Energy Security does not score downstream grid or upstream and midstream geothermal. Correlations are Spearman correlations across economies within a pair; the pooled row is the Pearson correlation of within-pair percentile ranks over all cells of the stage. A component is shared when the other pillar contains a component of the same category with identical observed values; its effective weight is its category weight times its within-category weight, averaged over cells. Non-shared pillars are the weighted means of the remaining components with the weights renormalised; they are undefined where every component of $\ES$ is shared, as in 7 downstream pairs, and the last two correlation rows use only the pairs in which both are defined. The covariance shares decompose the within-pair covariance of the pillar levels as in equation~\eqref{eq:joint_cov}, pooled over pairs with weights $n-1$: the shared term is $\operatorname{Cov}(A^{\ES},A^{\EO})$, the non-shared term $\operatorname{Cov}(B^{\ES},B^{\EO})$, and the two cross terms make up the remainder. Regions are defined by the within-pair medians of the two pillars, a cell at the median counting as low. Source: \texttt{index\_outputs.rds}; \texttt{scripts/paper/fig\_joint.R}.
\end{tablenotes}
\end{threeparttable}
\end{table}


\paragraph{Which categories carry the association.} Figure~\ref{fig:joint_corr} reports the
average within-pair Spearman correlation between every pair of categories that vary within
pairs. The cross-pillar block (bottom left) is the informative part. The three categories that
enter both pillars with identical components, Consumption, Energy Access and Investment, have a
correlation of one with their counterparts. Trade, which shares two components, has 0.62;
Production, built from different sources and transformations of the same physical quantity, has
0.27; and Foreign Dependency, whose shared components are missing for many economies and imputed
under different rules in the two pillars, has 0.04. Between categories with different names the
median cross-pillar correlation is 0.14 over the 47 category pairs that co-occur in at least one
technology-stage pair with 20 or more rankable economies; the largest is 0.76, between the
security pillar's Minerals Trade and the opportunity pillar's Trade in the five upstream mineral
pairs discussed above. The two categories that are distinctive to opportunity, Technology Demand
and Cost Competitiveness, are nearly orthogonal to every category of security: their absolute
correlations are at most 0.12 and 0.17, the most negative being $-0.17$, between Cost
Competitiveness and Reserves. The matrix also exposes a redundancy within $\ES$. Across the
seven upstream pairs in which both are scored, the average rank correlation between Production
and Reserves is 0.79. In the six pairs built on critical minerals (batteries, grid, electric
vehicles, green hydrogen, solar and wind) it ranges from 0.72 to 0.91: both are zero-anchored
indices over the technology's mineral basket (Sections~\ref{subsec:es_production}
and~\ref{subsec:es_reserves}), between 62 and 85 per cent of the rankable economies record zero
on both and tie at the bottom, and the ranks of the holders largely coincide, so that upstream
the security pillar counts mineral endowment twice over much of its range. Upstream nuclear,
scored on the identified resources and the mine production of the OECD-NEA and IAEA Red Book, has
the lowest value (0.61), because many holders of uranium resources do not mine them. Energy
Prices in both pillars and Technological Readiness in $\EO$ do not vary across economies within
any pair and so cannot affect within-pair rankings; they are omitted from the figure.

\begin{figure}[tbp]
  \centering
  \includegraphics[width=\textwidth]{joint_corr}
  \caption{Correlation among the categories of the two pillars. Each tile is the average, over
  the technology-stage pairs in which both categories are present and vary across economies (at
  least 20 rankable economies), of the Spearman correlation between the two category scores
  across economies within a pair; the number of pairs behind a tile ranges from two to 22.
  Categories are labelled by pillar; the black lines separate the block of cross-pillar
  correlations (bottom left) from the two within-pillar blocks. Numbers are printed for
  correlations of at least 0.5 in absolute value and, in bold, for categories that enter both
  pillars under the same name. A dash marks category pairs that never occur together in a
  technology-stage pair. Energy Prices (both pillars) and Technological Readiness ($\EO$) are
  constant within every pair and are omitted. Diverging scale centred at zero. Sample: the 3{,}157
  rankable cells of Figure~\ref{fig:joint_scatter}. Source: category scores in
  \texttt{index\_outputs.rds}; \texttt{scripts/paper/fig\_joint.R}.}
  \label{fig:joint_corr}
\end{figure}

\paragraph{Regions of the space.} Pooling all pairs, 15.1 per cent of cells combine above-median
opportunity with security at or below the median, 34.6 per cent are leaders on both, 13.2 per
cent combine above-median security with opportunity at or below the median, and 37.1 per cent
are low on both. Laggards outnumber leaders because cells tied at a median are counted as low,
and ties are common where a pillar rests on a zero-anchored component; counting tied cells as
high instead gives 12.6, 37.7, 15.3 and 34.4 per cent, and leaders then outnumber laggards. The
shares are largely the rank correlation restated. Under independence each region would hold a
quarter of the cells; a bivariate normal copula with each pair's observed Spearman correlation
implies an average off-diagonal share of 14.6 per cent, against 14.2 per cent observed, and
across pairs the observed share of opportunity without security tracks the implied one closely
(correlation 0.90). The off-diagonal regions are accordingly largest upstream, where the pillars
are least related (19.5 and 17.2 per cent), and smallest downstream (12.4 and 10.2 per cent).
Across technologies the share of cells with opportunity but not security ranges from 9.5 per
cent in electric vehicles to 20.2 per cent in green hydrogen and nuclear.

Among pairs (Figure~\ref{fig:joint_quadrants}), the share is highest in midstream green hydrogen
(45 per cent) and upstream nuclear (28 per cent), and in both it is largely an artefact of ties.
In midstream green hydrogen, 34 of the 40 rankable economies share the lowest value of
$\ES_{its}$, which is therefore also the median. In upstream nuclear, 87 of the 175 rankable
economies share the value 0.172, Energy Prices plus an imputed Trade category; all of them have
neither identified uranium resources nor mine production, and the common value is the pair's
median. A cell is thus above the security median only if the economy holds or mines uranium or
has an observed uranium trade that scores above the imputed value, and 37 of the 49 cells in the
region of opportunity without security are economies at the common value whose opportunity is
above the median. Counting tied cells as high empties the region in midstream green hydrogen and
lowers it to 6.9 per cent in upstream nuclear, and excluding the two pairs lowers the pooled
share from 15.1 to 14.0 per cent. The next pairs, upstream green hydrogen (24 per cent),
upstream coal (22 per cent) and downstream batteries (21 per cent), have few or no ties at the
security median, and their shares are near what the normal copula implies for their
correlations.

\begin{figure}[tbp]
  \centering
  \includegraphics[width=\textwidth]{joint_quadrants}
  \caption{Regions of the opportunity-security space, by technology and stage. Each bar divides
  the rankable economies of one technology-stage pair among the four regions of
  Figure~\ref{fig:joint_scatter}, which are defined by the within-pair medians of $\ES_{its}$ and
  $\EO_{its}$ over the rankable economies, a cell at the median counting as low. Panels are
  stages; within a panel, pairs are ordered by the share with opportunity but not security. The
  number to the right of a bar is the pair's count of rankable economies. Orange and blue are the
  colours of the Economic Opportunity and Energy Security pillars and mark the regions in which
  only that pillar is above its median; the darker gray marks leaders (both above) and the
  lighter gray laggards (neither). Ties at the median of $\ES_{its}$, which count as low, make the
  two longest orange segments: in midstream green hydrogen 34 of 40 economies share the lowest
  value of $\ES_{its}$, and in upstream nuclear 87 of 175 economies without identified uranium
  resources or mine production share the median value. The three pairs that Energy Security does
  not score (downstream grid, upstream and midstream geothermal) are absent. Source: as
  Figure~\ref{fig:joint_scatter}.}
  \label{fig:joint_quadrants}
\end{figure}

\paragraph{The posture of the coalition.} The coalition members sit mostly in the leaders'
region: 49.3 per cent of their 1{,}500 rankable cells are above both medians, 14.3 per cent have
security without opportunity, 23.6 per cent are low on both, and 12.8 per cent have opportunity
without security, below the 15.1 per cent of all rankable cells (Figure~\ref{fig:joint_members}).
The two tiers differ. The core members have 55.1 per cent of their cells among the leaders and
16.4 per cent with opportunity but not security, slightly above the share among all rankable
cells. The partners, many of which enter the coalition as major suppliers of upstream or
midstream inputs, have 40.4 per cent among the leaders, 21.4 per cent with security but not
opportunity and only 7.4 per cent with opportunity but not security. The concentration among
leaders is partly mechanical. Several components of both pillars rise with the scale of an
economy, among them world market shares, export size, investment and installed capacity, and the
core consists of economies with GDP of at least USD 100 billion, which therefore tend to sit above
both medians; the United States is a leader in 29 of its 30 rankable pairs. The median core
member has 16.7 per cent of its rankable pairs in the opportunity-without-security region, the
median partner 6.9 per cent.

The members with the most such pairs are European core economies whose security scores, which
reward domestic resources and production and low import exposure, are below the median in many
pairs while their opportunity scores are above it: Austria, Portugal and Luxembourg with 11 of 30
pairs each, Ireland with 10 and Croatia with 9. Their counts fall only to 9, 10, 10, 8 and 7 when
the two pairs dominated by ties are left out, so they do not rest on the tie artefact. Among
partners the largest counts are five, for Thailand and the Philippines (of 29 pairs each),
followed by four for Malaysia (of 29), Viet Nam and Qatar (of 26 each). Thirteen members have no
pair in the region, five in the core (Canada, Japan, the United Kingdom, Australia and Ukraine)
and eight partners (T\"urkiye, Brazil, Argentina, the Democratic Republic of the Congo, the
Dominican Republic, Gabon, Namibia and Ghana). The members with the most pairs of security
without opportunity are Brazil (19 of 29), Peru (17 of 28) and T\"urkiye (16 of 29): resource
holders whose reserves and production secure supply while their opportunity scores are below the
median. Brazil, which qualifies as a partner on its share of world mineral reserves alone
(Section~\ref{sec:coalition}), is the clearest case. Across members, the pairs most often in the
opportunity-without-security region are upstream gas, upstream oil and midstream green hydrogen
(14 members each), upstream coal (13), downstream solar (12), upstream nuclear (11) and upstream
grid (10). Two of these, upstream nuclear and midstream green hydrogen, are the tie artefacts
described above; without them the coalition's share in the region falls from 12.8 to 11.8 per
cent. The number of rankable pairs differs across members, from 22 to 30 in the core (median 30)
and from 9 to 29 among partners (median 27), so the shares of small partners such as Costa Rica,
Gabon and Namibia (9 pairs each), Guatemala (10), the Dominican Republic and Ghana (11 each) rest
on few cells.

\begin{figure}[tbp]
  \centering
  \includegraphics[width=\textwidth]{joint_members}
  \caption{The coalition in the opportunity-security space. Each bar divides one coalition
  member's rankable technology-stage pairs among the four regions of
  Figure~\ref{fig:joint_scatter}, which are defined by the within-pair medians of $\ES_{its}$ and
  $\EO_{its}$ over all rankable economies, a cell at the median counting as low. Left: the 31
  core members; right: the 26 partners (Section~\ref{sec:coalition}). Members are ordered within
  each tier by the share of their pairs with opportunity but not security; the fraction to the
  right of a bar is the count of such pairs over the member's rankable pairs, which number from 9
  to 30. The vertical line marks the share of all rankable cells in that region (15 per cent).
  Colours as in Figure~\ref{fig:joint_quadrants}. Upstream nuclear and midstream green hydrogen,
  in which many economies tie at the median of $\ES_{its}$, account for 25 of the 192 coalition
  pairs in the orange segments. Source: as Figure~\ref{fig:joint_scatter}.}
  \label{fig:joint_members}
\end{figure}

\paragraph{The within-country view.} The Strategic Index compares the sectors of one country, and
by equation~\eqref{eq:joint_var} what matters there is the joint distribution of the pillar
levels across those sectors. We compute it over the Strategic Index cells of each country. Of the
174 countries with a Strategic Index, 7 have fewer than five such cells (six of them one) and
are left out; the remaining 167 have between 5 and 30, with a median of 22, and 50 of them have
fewer than ten. The Pearson correlation of $\EO_{its}$ and $\ES_{its}$ averages 0.46 (median
0.48) and is positive for 164 of the 167 countries, ranging from $-0.56$ for the Maldives to 0.84
for Iran; the Spearman correlation averages 0.47, and the correlation of $\EO_{its}$ with the need
term $1-\ES_{its}$ is its negative (Spearman mean $-0.47$, from $-0.86$ for the Lao PDR to 0.70
for the Maldives). None of the three countries with a negative correlation (Afghanistan, Bhutan
and the Maldives) is a coalition member. Among the 57 coalition members the Pearson correlation
averages 0.52, from 0.27 for New Zealand to 0.77 for Viet Nam. Security varies more than
opportunity across a country's sectors (mean standard deviations of 0.15 and 0.07), and the
covariance term of equation~\eqref{eq:joint_var} removes on average 36 per cent of
$\operatorname{Var}_i(\EO)+\operatorname{Var}_i(\ES)$ from the variance of the gap; for coalition
members the average is 42 per cent, from 16.5 per cent for Ghana to 66 per cent for India. The
standard deviation of $\Delta_{its}$ across a country's sectors is accordingly 0.13 on average
(from 0.068 for the Lao PDR to 0.26 for the Gambia), so that the pillar part of the index,
$0.3\,\Delta_{its}$, varies with a standard deviation of about 0.04. This compression is not
simply the differences between technology-stage pairs that are common to every country: the
pairs' mean values of the two pillars are only weakly related across the 30 pairs (Pearson
correlation 0.33), and after the pair means over all rankable economies are subtracted the
within-country correlation rises to an average of 0.67 (median 0.73). Subtracting the pair means
lowers the average standard deviation of $\Delta_{its}$ from 0.13 to 0.08, so part of the spread
of the gap across a country's sectors is the gap between the pillars' pair means, which is the
same for every country.

\paragraph{The regions and the Strategic Index.} The Strategic Index of equation~\eqref{eq:si}
adds $0.3\,\EO_{its}$ to $0.3\,(1-\ES_{its})$ and so rewards opportunity without security. Every
one of its 3{,}154 cells, for 174 countries, lies in the sample of this section, since the index
requires rankable values of both pillars. On average it does favour the region: the mean index is
0.563 there, against 0.545 among laggards, 0.522 among leaders and 0.511 where security exceeds
opportunity. A country's highest-ranked sector, over all stages and before any eligibility rule,
lies in the region for 47 of the 174 countries (27.0 per cent; 15 of the 57 coalition members),
whereas 15.2 per cent of all index cells do. Much of this is upstream nuclear, the top sector of
62 countries, 43 of which share the common value of $\ES_{its}$ described above: that value gives
them a large need term whichever side of the opportunity median they fall, and the median split
places 23 of the 62 in the region and 30 among the laggards. The endowment screen of
Section~\ref{subsec:si_allied}, which removes upstream extraction stages without a recorded
resource base from the eligible set, leaves 169 countries with at least one cell it does not flag,
and for 35 of them (20.7 per cent) the highest-ranked such cell lies in the region.

The allied picks apply the screen and also confine the choice to upstream and midstream sectors
other than geothermal and green hydrogen. All 57 members have one, 43 of them midstream. Eight
lie in the region of opportunity without security: Austria, Belgium, Italy, the Netherlands,
Romania and Viet Nam in midstream batteries, and Indonesia and Slovakia in midstream nuclear. Of
the rest, 21 are among the leaders, 14 have security without opportunity and 14 are among the
laggards. The eight are 14.0 per cent of the picks, against 6.7 per cent of the 646 eligible
sectors of the same members, of which 61.0 per cent are leaders, so the picks lean towards the
region without being confined to it. None of the eight upstream nuclear picks lies in the region:
all eight members hold or mine uranium, three are leaders and five have security without
opportunity. Four of the nine midstream nuclear picks come from small partners with few eligible
sectors: Costa Rica, the Dominican Republic and Guatemala have four each, the midstream stages of
nuclear, oil, gas and coal, and Namibia five, the same four and upstream nuclear, in which its
uranium output gives it a high security score and so a small need. This is not an effect of a
security floor. In midstream nuclear $\ES_{its}$ is built from the Trade category alone and is
compressed around its median (quartiles 0.435, 0.482 and 0.530), so the four partners' security
ranks of 0.50 to 0.71 place them at or above it (Costa Rica among the laggards, the other
three with security without opportunity), and among the midstream candidates the technology
weight $g_t$ of equation~\eqref{eq:si}, 0.999 for nuclear against 0.422 for gas, 0.024 for coal
and zero for oil, favours nuclear. The index is thus not organised by the regions. It combines
pillar levels, which carry the large differences across technology-stage pairs, with the stage
and technology weights, whereas the regions are defined by positions within pairs: across a
country's sectors the index has a median Spearman correlation of 0.63 with the level gap
$\EO_{its}-\ES_{its}$ but of only 0.19 with the gap between the two within-pair percentile ranks.
A country's sector ranking under the Strategic Index therefore reflects which technology-stage
pairs have high opportunity and low security scores for everyone at least as much as where the
country stands relative to other economies; Section~\ref{sec:strategic} discusses this property
of the index.

\paragraph{Caveats.} Five features of the construction qualify these results. First, percentile
ranks discard distances, and the median split turns small differences near the median into
differences of region; the shares should be read as descriptions of rank positions, not of
magnitudes. Where many economies tie at a median, the convention that a tied cell counts as low
decides the region of the whole tie group. This matters mainly upstream, through the common
value of upstream nuclear: counting ties as high lowers the upstream share of opportunity
without security from 19.5 to 15.5 per cent and the coalition's share from 12.8 to 10.8 per cent,
and it leaves the ordering of the stages unchanged. Second, the regions are defined relative to
the rankable economies of each pair, from 40 to 175 of them. The rankability rule removes
territories and economies whose scores rest mostly on imputation, so the comparison set leans
towards economies with good statistical coverage, and ``below the median'' means below the median
of that set. Third, the share of the association that is mechanical is measured on ranks only in
the 23 pairs in which both pillars retain a component of their own; in the other seven the
question does not arise, because security there is a subset of opportunity. The covariance
decomposition covers all 30 pairs but is a statement about levels, and its pooled shares are
dominated by the pairs with the largest covariances. Fourth, the within-country statistics are
computed over at most 30 sectors whose pillar values are weighted means over category sets that
differ across pairs (Section~\ref{subsec:norm_pillar}), so they compare bundles that are not
identical, and for the 50 countries with fewer than ten Strategic Index cells they rest on few
observations. Fifth, the endowment screen is a rule for the allied picks, not part of either
pillar or of the index; the regions and correlations above are computed without it. More
generally, the overlap of inputs is the familiar difficulty of composite indices that share
components, whose implied trade-offs are hard to see in the aggregate
\citep{ravallion2012mashup}; inspecting the correlation structure among indicators, as
recommended by \citet{oecd2008handbook}, is what reveals it here.
